% Options for packages loaded elsewhere
\PassOptionsToPackage{unicode}{hyperref}
\PassOptionsToPackage{hyphens}{url}
\documentclass[
  a4paper,
]{ltjsarticle}
\usepackage{xcolor}
\usepackage[margin=25mm]{geometry}
\usepackage{amsmath,amssymb}
\setcounter{secnumdepth}{-\maxdimen} % remove section numbering
\usepackage{iftex}
\ifPDFTeX
  \usepackage[T1]{fontenc}
  \usepackage[utf8]{inputenc}
  \usepackage{textcomp} % provide euro and other symbols
\else % if luatex or xetex
  \usepackage{unicode-math} % this also loads fontspec
  \defaultfontfeatures{Scale=MatchLowercase}
  \defaultfontfeatures[\rmfamily]{Ligatures=TeX,Scale=1}
\fi
\usepackage{lmodern}
\ifPDFTeX\else
  % xetex/luatex font selection
\fi
% Use upquote if available, for straight quotes in verbatim environments
\IfFileExists{upquote.sty}{\usepackage{upquote}}{}
\IfFileExists{microtype.sty}{% use microtype if available
  \usepackage[]{microtype}
  \UseMicrotypeSet[protrusion]{basicmath} % disable protrusion for tt fonts
}{}
\makeatletter
\@ifundefined{KOMAClassName}{% if non-KOMA class
  \IfFileExists{parskip.sty}{%
    \usepackage{parskip}
  }{% else
    \setlength{\parindent}{0pt}
    \setlength{\parskip}{6pt plus 2pt minus 1pt}}
}{% if KOMA class
  \KOMAoptions{parskip=half}}
\makeatother
\usepackage{longtable,booktabs,array}
\newcounter{none} % for unnumbered tables
\usepackage{calc} % for calculating minipage widths
% Correct order of tables after \paragraph or \subparagraph
\usepackage{etoolbox}
\makeatletter
\patchcmd\longtable{\par}{\if@noskipsec\mbox{}\fi\par}{}{}
\makeatother
% Allow footnotes in longtable head/foot
\IfFileExists{footnotehyper.sty}{\usepackage{footnotehyper}}{\usepackage{footnote}}
\makesavenoteenv{longtable}
\usepackage{graphicx}
\makeatletter
\newsavebox\pandoc@box
\newcommand*\pandocbounded[1]{% scales image to fit in text height/width
  \sbox\pandoc@box{#1}%
  \Gscale@div\@tempa{\textheight}{\dimexpr\ht\pandoc@box+\dp\pandoc@box\relax}%
  \Gscale@div\@tempb{\linewidth}{\wd\pandoc@box}%
  \ifdim\@tempb\p@<\@tempa\p@\let\@tempa\@tempb\fi% select the smaller of both
  \ifdim\@tempa\p@<\p@\scalebox{\@tempa}{\usebox\pandoc@box}%
  \else\usebox{\pandoc@box}%
  \fi%
}
% Set default figure placement to htbp
\def\fps@figure{htbp}
\makeatother
\setlength{\emergencystretch}{3em} % prevent overfull lines
\providecommand{\tightlist}{%
  \setlength{\itemsep}{0pt}\setlength{\parskip}{0pt}}
% Glyph map for lualatex (newunicodechar). Maps symbols that may lack font glyphs.
\usepackage{newunicodechar}
\usepackage{amssymb}
\newunicodechar{∎}{\ensuremath{\blacksquare}}
\newunicodechar{✓}{\ensuremath{\checkmark}}
\newunicodechar{✗}{\ensuremath{\times}}
\newunicodechar{⟹}{\ensuremath{\Longrightarrow}}
\newunicodechar{⟺}{\ensuremath{\Longleftrightarrow}}
\newunicodechar{↪}{\ensuremath{\hookrightarrow}}
\newunicodechar{⌊}{\ensuremath{\lfloor}}
\newunicodechar{⌋}{\ensuremath{\rfloor}}
\newunicodechar{≃}{\ensuremath{\simeq}}
\newunicodechar{≈}{\ensuremath{\approx}}
\newunicodechar{≤}{\ensuremath{\leq}}
\newunicodechar{≥}{\ensuremath{\geq}}
\newunicodechar{≠}{\ensuremath{\neq}}
\newunicodechar{→}{\ensuremath{\rightarrow}}
\newunicodechar{←}{\ensuremath{\leftarrow}}
\newunicodechar{↔}{\ensuremath{\leftrightarrow}}
\newunicodechar{⊥}{\ensuremath{\perp}}
\newunicodechar{√}{\ensuremath{\surd}}
\newunicodechar{∑}{\ensuremath{\sum}}
\newunicodechar{∏}{\ensuremath{\prod}}
\newunicodechar{∞}{\ensuremath{\infty}}
\newunicodechar{∈}{\ensuremath{\in}}
\newunicodechar{∀}{\ensuremath{\forall}}
\newunicodechar{∣}{\ensuremath{\mid}}
\newunicodechar{γ}{\ensuremath{\gamma}}
\newunicodechar{λ}{\ensuremath{\lambda}}
\newunicodechar{μ}{\ensuremath{\mu}}
\newunicodechar{ρ}{\ensuremath{\rho}}
\newunicodechar{σ}{\ensuremath{\sigma}}
\newunicodechar{φ}{\ensuremath{\varphi}}
\newunicodechar{ψ}{\ensuremath{\psi}}
\newunicodechar{θ}{\ensuremath{\theta}}
\newunicodechar{Δ}{\ensuremath{\Delta}}
\newunicodechar{Π}{\ensuremath{\Pi}}
\newunicodechar{Λ}{\ensuremath{\Lambda}}
\newunicodechar{τ}{\ensuremath{\tau}}
\newunicodechar{ε}{\ensuremath{\varepsilon}}
\newunicodechar{ω}{\ensuremath{\omega}}
\newunicodechar{π}{\ensuremath{\pi}}
\newunicodechar{ν}{\ensuremath{\nu}}
\newunicodechar{±}{\ensuremath{\pm}}
\newunicodechar{×}{\ensuremath{\times}}
\newunicodechar{·}{\ensuremath{\cdot}}
\usepackage{bookmark}
\IfFileExists{xurl.sty}{\usepackage{xurl}}{} % add URL line breaks if available
\urlstyle{same}
\hypersetup{
  pdftitle={Analytic Exact Solution of the 90-Degree High-Symmetry Floor --- The \textbackslash sigma\_\{\textbackslash max\} Eigenmode of the Integer Antisymmetric Generator Fixed by Z4 Phase Labels{,} and Its Closed-Form Spectrum},
  pdfauthor={Noriaki Kihara (WF System Co., Ltd.), ORCID 0009-0004-6753-4020},
  hidelinks,
  pdfcreator={LaTeX via pandoc}}

\title{Analytic Exact Solution of the 90-Degree High-Symmetry Floor ---
The \(\sigma_{\max}\) Eigenmode of the Integer Antisymmetric Generator
Fixed by Z4 Phase Labels, and Its Closed-Form Spectrum}
\author{Noriaki Kihara (WF System Co., Ltd.), ORCID 0009-0004-6753-4020}
\date{2026-09-12}

\begin{document}
\maketitle

\textbf{Series:} Foundations of Self-Consistent Relational-Wave Closed
Systems (Paper 3 of 10)\\
\textbf{Author:} Noriaki Kihara (WF System Co., Ltd.)　\textbf{Date:}
2026-09-12\\
\textbf{Version DOI:} 10.5281/zenodo.22728996\\
\textbf{Concept DOI:} 10.5281/zenodo.22728995\\
\textbf{ORCID:} 0009-0004-6753-4020\\
\textbf{Zenodo:} https://zenodo.org/records/22728996\\

\begin{center}\rule{0.5\linewidth}{0.5pt}\end{center}

\subsection{Abstract}\label{abstract}

The 90-degree high-symmetry floor, established in Paper 2 as satisfying
the ignition qualification, is here constructed \textbf{not} by
iterative approximation (make\_parent, \texttt{iters=1200},
\texttt{tol=1e-12}) but \textbf{analytically and exactly as an integer
eigenvalue problem}, and its closed form is derived for
\textbf{arbitrary \(N\)}. Because the floor phases, up to a global
phase, are quantized to \(\{0,90,180,270\}^\circ\), assigning each edge
a Z4 label \(k_e\in\{0,1,2,3\}\) makes the generator
\(K_{ef}=A_{ef}\sin(90^\circ(k_f-k_e))\) an \textbf{integer matrix
(entries \(0,\pm1\))}. For the deterministic seed \(k_e=(i+j)\bmod4\),
what determines the generator is not mod 4 but only the \textbf{parity}
\(p_e=(i+j)\bmod2\), and the real symmetric matrix
\(W_{ef}=A_{ef}\sin^2\psi_{ef}\) of Paper 1 becomes a bipartite graph
\(W_{ef}=A_{ef}\mathbf 1_{p_e\neq p_f}\) (\(L(K_N)\) split by the parity
of the sum of vertex indices). Therefore the floor can be constructed
\textbf{in one shot} by \(k_e=(i+j)\bmod4\to W\to\) Perron vector
\(r>0\to z=Dr\) (\(D=\operatorname{diag}(e^{i90^\circ k_e})\), Paper 1
\(D^{-1}KD=iW\)), and \textbf{discrete-label iteration is essentially
unnecessary}. The floor is the eigenvector of the \textbf{negative
eigenvalue of largest absolute value} \(\mu=-\sigma_{\max}\) of \(iK\),
and its phase advances as \(d\varphi/d\tau=-\mu=\sigma_{\max}>0\),
consistent with the phase equation. From the reduced Perron problem of
the bipartite structure (\(2\times2\) for even \(N\), \(3\times3\) for
odd \(N\)), \(\sigma^2\) (even \(N(N-2)\) / odd \(N^2-2N-1\)), per-edge
squared amplitude (rational within each parity class), and the full
spectrum (\(BB^{\mathsf T}\)) all come out in closed form for arbitrary
\(N\) (\textbf{verified symbolically in \(N\)-general with sympy, and
numerically matched over \(N=3\)--\(40\)}). Moreover, for general \(N\)
the number of nonzero spectral values is \(O(N)\), and for even
\(N\ge6\) / odd \(N\ge7\) it is exactly \(2(N-1)\)
\(=\operatorname{rank}K=2(N-1)\) with kernel dimension \((N-1)(N-4)/2\)
(\(N=3,4,5\) are degenerate exceptions): of the
\(M=N(N-1)/2=O(N^2)\)-dimensional relational-wave space, only \(O(N)\)
dimensions are active. Over all \(N=3,\ldots,40\) (38 floors), the
one\_step relative-equilibrium residual is
\(\sim10^{-15}\)--\(1.2\times10^{-14}\) (better than make\_parent's
\(\sim10^{-13}\)), the phase lies exactly on the axes, and
\(\|z\|^2=1\). \textbf{Zero-closure \(\sum z^2=0\) is not an extra
condition imposed on the floor but a general theorem that follows
automatically from the mode being a bipartite eigenmode
(\(\|x\|^2=\|y\|^2\))} (holding not only for the Perron mode but for
every nonzero eigenmode), and the connectivity of the bipartite \(W\) is
also proven for arbitrary \(N\ge3\) (recovering the applicability
condition of the Perron argument in Paper 1). The floor can be written
down completely algebraically as ``\(\sigma=\sqrt{\text{integer}}\),
per-edge amplitude \(=\sqrt{\text{rational}}\) (fixed by the parity
class), phase \(\in\{0,90,180,270\}\).'' Classification: analytic exact
solution (parity bipartite structure + reduced Perron problem) +
symbolic (arbitrary \(N\)) / numerical (\(N=3\)--\(40\)) verification.

\begin{center}\rule{0.5\linewidth}{0.5pt}\end{center}

\subsection{1. The Problem --- Making the Approximate Floor
Exact}\label{the-problem-making-the-approximate-floor-exact}

The make\_parent floor used in Paper 2 is a solution that iteratively
approximates the self-consistent eigenmode, with residual
\(\sim10^{-13}\). Merely copying an approximation stopped by an
iteration cap, tolerance cut, and damping is meaningless. The 90-degree
floor has the structural property that its phases are quantized to a
discrete set (Z4), and using this it can be solved \textbf{algebraically
and exactly as an integer eigenvalue problem}. This paper constructs
that exact solution for all \(N\) (including odd \(N\)) and fixes the
closed form.

\subsection{2. Principle --- From Z4 Labels to the Integer Antisymmetric
Generator}\label{principle-from-z4-labels-to-the-integer-antisymmetric-generator}

Quantize the floor phases, up to a global phase, into the four quadrants
and assign each edge \(k_e\in\{0,1,2,3\}\). Since the phase difference
is \(90^\circ(k_f-k_e)\),

\[K_{ef}=A_{ef}\sin\!\big(90^\circ(k_f-k_e)\big),\qquad (k_f-k_e)\bmod4:\ 0\to0,\ 1\to+1,\ 2\to0,\ 3\to-1,\]

that is, \(K\) is an \textbf{integer antisymmetric matrix} with entries
\(0,\pm1\). The floor \(z\) is its \(\sigma_{\max}\) eigenmode, i.e.~the
eigenvector of the \textbf{negative eigenvalue of largest absolute
value} \(\mu=-\sigma_{\max}\) of \(iK\) (of the form
\(z=\sigma_{\max}u-iKu\), where \(u\) is the real eigenvector of
\(K^2u=-\sigma_{\max}^2u\)). \(iK\) is Hermitian with eigenvalues in
conjugate pairs \(\pm\sigma_j\), and \(\sigma_{\max}\) is the largest
singular value (spectral radius). The floor takes the branch
\(\mu=-\sigma_{\max}\) where the phase advances, and
\(d\varphi/d\tau=-\mu=\sigma_{\max}>0\) is consistent with the phase
equation \(r\dot\varphi=\sum r_f\sin^2\Delta\varphi\ge0\) (confirmed
with real data and code: for N=6 the make\_parent floor has
\(\mu=-4.868\), the integer-K floor \(\mu=-4.899\), both the smallest
eigenvalue). What the square roots of the eigenvalues of
\(K^{\mathsf T}K\) fix are the \textbf{singular values
\(\sigma=\sqrt{\text{integer}}\)}, whereas each per-edge amplitude
\(r_e\) is a component of the eigenvector belonging to it (= the Perron
vector of \(W\)), and in this construction it takes the closed form
\(r_e=\sqrt{\text{rational}}\) (§5).

Note that the complex self-consistent problem \(Kz=i\omega z\) is
exactly equivalent to the eigenvalue problem \(Wr=\omega r\) of the real
symmetric non-negative integer matrix \(W_{ef}=A_{ef}\sin^2\psi_{ef}\)
(\(D^{-1}KD=iW\), \(D=\operatorname{diag}(e^{i\varphi_e})\)), and the
existence and uniqueness of a positive-amplitude solution by
Perron--Frobenius when \(W\) is irreducible is \textbf{given in general
by Paper 1}. This paper, with that \(W\) = the concrete form of the
integer \(K\), solves \(\sigma=\rho(W)\) and each per-edge amplitude in
closed form (§5) (the Perron eigenvector of \(W\) and the
\(\sigma_{\max}\) mode of the integer \(K\) are the same floor;
confirmed \(r>0\) and agreement of the principal eigenvector for all
\(N\)).

\textbf{Bipartite structure from parity.} What determines \(W\) is not
mod 4 but only \textbf{mod 2} (parity) \(p_e=(i+j)\bmod2\). Since
\(\sin^2[90^\circ(k_f-k_e)]\) is \(1\) when \(k_f-k_e\) is odd
(\(p_e\neq p_f\)) and \(0\) when even (\(p_e=p_f\)),

\[W_{ef}=A_{ef}\,\mathbf 1_{p_e\neq p_f}.\]

That is, \(W\) is the \textbf{adjacency matrix of the bipartite graph}
obtained by splitting the vertices of \(L(K_N)\) (= the edges
\(\{i,j\}\) of \(K_N\)) by the parity \(p_e\) of the sum of vertex
indices. This \(W\) is \textbf{connected (irreducible)} for arbitrary
\(N\ge3\) (proof in §8: for the even-vertex set \(U\)
(\(|U|=\lceil N/2\rceil\ge2\)) and odd-vertex set \(V\), cross-parity
edges are bridged by internal edges, and same-parity edges are connected
via cross-parity edges). Hence the Perron vector \(r>0\) is unique
(Paper 1), and \(z=Dr\) (\(D=\operatorname{diag}(e^{i90^\circ k_e})\))
exactly satisfies \(Kz=i\sigma_{\max}z\). Therefore the floor can be
constructed in one shot from the seed labels by \(k\to W\to r\to z=Dr\),
and the discrete-label iteration (§3) is unnecessary for constructing
the floor (confirmed for all \(N=3\)--\(40\) that the \(W\) of the
iterative floor agrees with the seed-parity \(W\) to machine precision,
\(\|\text{diff}\|\le3\times10^{-27}\),
\texttt{parity二部構造\_解析的証明\_20260912/}). From this bipartite
structure the closed forms of §5 and §6 come out for arbitrary \(N\).

\textbf{Square-closure \(\sum_e z_e^2=0\) is a consequence of the
bipartite structure (not an extra condition).} For the bipartite
\(W=\begin{pmatrix}0&B\\B^{\mathsf T}&0\end{pmatrix}\), a nonzero
eigenmode \(r=(x,y)\) satisfies
\(By=\sigma x,\ B^{\mathsf T}x=\sigma y\), so
\(\sigma\|x\|^2=x^{\mathsf T}By=(B^{\mathsf T}x)^{\mathsf T}y=\sigma\|y\|^2\),
and since \(\sigma\neq0\), \(\|x\|^2=\|y\|^2\) (each \(1/2\) under
normalization). On the other hand, from \(z_e=e^{i90^\circ k_e}r_e\) we
have \(z_e^2=(-1)^{k_e}r_e^2=(-1)^{p_e}r_e^2\) (real), so

\[\sum_e z_e^2=\sum_{p_e=0}r_e^2-\sum_{p_e=1}r_e^2=\|x\|^2-\|y\|^2=0.\]

That is, zero-closure \(\sum z^2=0\) is not an extra condition imposed
on the floor but \textbf{arises automatically from being a bipartite
eigenmode} (holding not only for the Perron floor but for any nonzero
eigenmode of the bipartite \(W\); the real part and the cross-term
imaginary part vanish separately as a consequence of \(z_e^2\) being
real). Confirmed \(|\sum z^2|\le10^{-15}\) and
\(|\|x\|^2-\|y\|^2|\le10^{-15}\) for every nonzero eigenmode over all
\(N=3\)--\(40\).

\subsection{3. Construction of the Floor --- One Shot from the Seed
Labels via Perron (Iteration for Confirmation
Only)}\label{construction-of-the-floor-one-shot-from-the-seed-labels-via-perron-iteration-for-confirmation-only}

By §2 the floor can be constructed \textbf{directly without iteration}
from the seed labels \(k_e=(i+j)\bmod4\) (edge \(e=\{i,j\}\),
\(0\le i<j\le N-1\)):

\begin{enumerate}
\def\labelenumi{\arabic{enumi}.}
\tightlist
\item
  Build the bipartite matrix \(W_{ef}=A_{ef}\mathbf 1_{p_e\neq p_f}\)
  using the parity \(p_e=(i+j)\bmod2\).
\item
  Take the Perron vector \(r>0\) of \(W\) and \(\sigma_{\max}=\rho(W)\)
  (\(W\) connected, Paper 1).
\item
  With \(D=\operatorname{diag}(e^{i90^\circ k_e})\), the floor is
  \(z=Dr\). By Paper 1 \(D^{-1}KD=iW\), so \(Kz=i\sigma_{\max}z\)
  (relative equilibrium).
\end{enumerate}

We confirmed for all \(N=3\)--\(40\) that this direct construction
\(z=Dr\) satisfies a one\_step relative-equilibrium residual
\(\le1.2\times10^{-14}\). The conventional generator
\texttt{make\_90deg\_floor\_analytic\_v1.py} obtains the floor from the
seed via discrete-label iteration (re-quantize the eigenvector of the
eigenvalue of largest absolute value \(\mu=-\sigma_{\max}\) of \(iK\)
into the four quadrants → best-match to the previous labels over the
four global-phase choices → stop at a fixed point / gauge-equivalent
cycle, adopting the minimum-residual floor), but the \textbf{\(W\) of
the iterative floor agrees with the seed-parity \(W\) to machine
precision} (\(\|\text{diff}\|\le3\times10^{-27}\)), and the floor is
identical to the direct seed construction up to a \textbf{sign gauge}
(\(z_e\to-z_e\), \(k_e\to k_e+2\), \(K\to SKS\) with
\(Kz=i\sigma_{\max}z\) and \(W\) invariant) (after gauge alignment
\(\|\text{diff}\|\le8\times10^{-15}\)). Therefore the iteration is
unnecessary for constructing the floor and amounts to confirming that
the seed is a fixed point (up to the sign gauge). adjacency/one\_step
are extracted from the canonical source
\texttt{run\_N3\_N40\_stage123\_v1.py} (SHA-checked) (no physics
changed). For each floor's verification
(\texttt{build\_analytic\_floor}), the following are recorded:
\(\|z\|^2\), \(|\sum z|\), \(\Re(a^{\mathsf T}a-b^{\mathsf T}b)\),
\(a^{\mathsf T}b\), \(|\sum z^2|\), the one\_step relative-equilibrium
residual, the on-axis test, the number of quadrants, the number of
amplitude classes, and \(\sigma\).

\subsection{\texorpdfstring{4. Results (All \(N=3,\ldots,40\) · 38
Floors)}{4. Results (All N=3,\textbackslash ldots,40 · 38 Floors)}}\label{results-all-n3ldots40-38-floors}

For all floors the phase is exactly on-axis (\(0/90/180/270\)), closure
\(\sum z^2=0\) (both the real part and the cross-term imaginary part
machine-zero; a consequence of the bipartite structure of §2),
\(\|z\|^2=1\), and the relative-equilibrium residual is
\(\sim10^{-15}\)--\(1.2\times10^{-14}\).

\begin{itemize}
\tightlist
\item
  \textbf{\(N=3\)}: 3 quadrants (T-shape), \(|\sum z|=0.707\),
  \(\sigma=\sqrt2=1.414\).
\item
  \textbf{\(4\mid N\) (4,8,\ldots,40)}: centroid
  \(|\sum z|\sim10^{-15}\) (exactly zero), 2-valued amplitude.
\item
  \textbf{\(N\equiv2\ (\mathrm{mod}\,4)\) (6,10,\ldots,38)}: centroid
  \(0.236\to0.037\) (decreasing as \(N\) increases), 2-valued amplitude.
\item
  \textbf{Odd \(N\) (5,7,\ldots,39)}: centroid \(0.154\to0.019\)
  (decreasing), 3-valued amplitude. \textbf{Odd \(N\) are all floors
  too} (on-axis, zero closure).
\end{itemize}

An exactly-zero centroid occurs only for \(4\mid N\) (where complete
\(\pm\) symmetry of the quadrants is possible); otherwise it is a
residual centroid determined by parity imbalance and the Perron solution
(a determined quantity, not a minimization; closed form in §8; an exact
zero cannot be attained). Precision comparison (same \(N\)):

{\def\LTcaptype{none} % do not increment counter
\begin{longtable}[]{@{}lll@{}}
\toprule\noalign{}
\(N\) & make\_parent residual & integer-K analytic residual \\
\midrule\noalign{}
\endhead
\bottomrule\noalign{}
\endlastfoot
6 & \(2.3\times10^{-13}\) & \(5.4\times10^{-16}\) \\
7 & \(3.5\times10^{-14}\) & \(1.7\times10^{-15}\) \\
8 & \(8.2\times10^{-14}\) & \(5.7\times10^{-16}\) \\
9 & \(8.2\times10^{-14}\) & \(1.0\times10^{-15}\) \\
\end{longtable}
}

All 38 floors are collected in a complex-plane grid figure (Figure 1).

\begin{figure}
\centering
\pandocbounded{\includegraphics[keepaspectratio,alt={Figure 1: 90-degree theoretical floors (integer-K σ\_max analytic solution, all N=3..40). Phase 0/90/180/270 · exact relative equilibrium (residual \textasciitilde1e-15) · centroid exactly zero for 4\textbar N.}]{figs/p3_analytic_floor_en_fig1.png}}
\caption{Figure 1: 90-degree theoretical floors (integer-K σ\_max
analytic solution, all N=3..40). Phase 0/90/180/270 · exact relative
equilibrium (residual \textasciitilde1e-15) · centroid exactly zero for
4\textbar N.}
\end{figure}

\textbf{Figure 1} 90-degree theoretical floors (integer-K
\(\sigma_{\max}\) analytic solution, \(N=3\)--\(40\)).

\subsection{\texorpdfstring{5. Closed Form (From the Parity-Reduced
Perron Problem, Arbitrary
\(N\))}{5. Closed Form (From the Parity-Reduced Perron Problem, Arbitrary N)}}\label{closed-form-from-the-parity-reduced-perron-problem-arbitrary-n}

From the bipartite structure of §2, the Perron vector \(r\) takes a
constant value within each parity class. Counting the number of
adjacencies for each class yields a reduced Perron problem, and
\(\sigma\) and each per-edge amplitude are solved in closed form for
\textbf{arbitrary \(N\)} (sympy symbolic verification,
\texttt{parity二部構造\_解析的証明\_20260912/results/実行ログ}). Since
the integer \(K\) is antisymmetric, \(K^2\) is integer symmetric
(eigenvalues \(-\sigma^2\)), so \(\sigma=\sqrt{\text{integer}}\) and
each per-edge squared amplitude is rational.

\textbf{Even \(N=2m\) (2 classes).} Cross-parity edges (\(p=1\), \(i+j\)
odd, count \(N^2/4\)) and same-parity internal edges (\(p=0\), \(i+j\)
even, count \(N(N-2)/4\)). Since each cross-parity edge has \(N-2\)
same-parity edges adjacent and each same-parity edge has \(N\)
cross-parity edges adjacent, the class representative values \(x\)
(cross-parity) and \(y\) (same-parity) satisfy

\[\sigma x=(N-2)y,\qquad \sigma y=Nx\ \Longrightarrow\ \sigma^2=N(N-2),\]

and normalizing by \(\|z\|^2=1\) gives \(x^2=2/N^2\),
\(y^2=2/(N(N-2))\).

\textbf{Odd \(N\) (3 classes).} \(a=\tfrac{N+1}2\) even vertices,
\(b=\tfrac{N-1}2\) odd vertices. Cross-parity (\(ab=\tfrac{N^2-1}4\)),
even-side internal (\(\binom a2=\tfrac{N^2-1}8\)), odd-side internal
(\(\binom b2=\tfrac{(N-1)(N-3)}8\)). The representative values \(x\)
(cross-parity), \(y\) (even-side internal), \(z\) (odd-side internal)
satisfy

\[\sigma x=\tfrac{N-1}2y+\tfrac{N-3}2z,\quad \sigma y=(N-1)x,\quad \sigma z=(N+1)x\ \Longrightarrow\ \sigma^2=N^2-2N-1,\]

and normalizing gives \(x^2=\tfrac{2}{(N-1)(N+1)}\),
\(y^2=\tfrac{2(N-1)}{(N+1)(N^2-2N-1)}\),
\(z^2=\tfrac{2(N+1)}{(N-1)(N^2-2N-1)}\).

\textbf{\(\sigma^2\)}: even \(N(N-2)\) / odd \(N^2-2N-1\). Sequence
\(N=3\)--\(20\):
\(2,8,14,24,34,48,62,80,98,120,142,168,194,224,254,288,322,360\).

\textbf{Squared amplitude (normalized \(\|z\|^2=1\), value × count) ---
the class assignment is fixed by parity.}

Even \(N\) (2 classes, each class total power \(1/2\)):

{\def\LTcaptype{none} % do not increment counter
\begin{longtable}[]{@{}lll@{}}
\toprule\noalign{}
Squared amplitude & Class (edge) & Count \\
\midrule\noalign{}
\endhead
\bottomrule\noalign{}
\endlastfoot
\(2/N^2\) & cross-parity (\(i+j\) odd) & \((N/2)^2\) \\
\(2/(N(N-2))\) & same-parity internal (\(i+j\) even) & \(N(N-2)/4\) \\
\end{longtable}
}

Odd \(N\) (3 classes; \(N=3\) has 0 odd-side internal edges = T-shape):

{\def\LTcaptype{none} % do not increment counter
\begin{longtable}[]{@{}
  >{\raggedright\arraybackslash}p{(\linewidth - 4\tabcolsep) * \real{0.3333}}
  >{\raggedright\arraybackslash}p{(\linewidth - 4\tabcolsep) * \real{0.3333}}
  >{\raggedright\arraybackslash}p{(\linewidth - 4\tabcolsep) * \real{0.3333}}@{}}
\toprule\noalign{}
\begin{minipage}[b]{\linewidth}\raggedright
Squared amplitude
\end{minipage} & \begin{minipage}[b]{\linewidth}\raggedright
Class (edge)
\end{minipage} & \begin{minipage}[b]{\linewidth}\raggedright
Count
\end{minipage} \\
\midrule\noalign{}
\endhead
\bottomrule\noalign{}
\endlastfoot
\(2/((N-1)(N+1))\) & cross-parity (\(i+j\) odd) & \((N^2-1)/4\) \\
\(2(N-1)/((N+1)(N^2-2N-1))\) & even-side internal (both \(i,j\) even) &
\((N^2-1)/8\) \\
\(2(N+1)/((N-1)(N^2-2N-1))\) & odd-side internal (both \(i,j\) odd) &
\((N-1)(N-3)/8\) \\
\end{longtable}
}

In each case the total of the counts is \(M=N(N-1)/2\). The closed form
is verified symbolically with sympy (arbitrary \(N\)) and matches the
numerical per-edge squared amplitudes over \(N=3\)--\(40\) (error
\(\le1.2\times10^{-16}\)). The centroid is exactly zero for even \(N\)
because the 2 classes, each with \(1/2\) power, cancel completely as
\(\pm\). For odd \(N\) the \(\pm\) cancellation is incomplete across the
3 classes, leaving a small residual centroid (closed form in §8).

\subsection{\texorpdfstring{6. Eigen-Oscillation Spectrum and Active
Dimension
(\(O(N)\))}{6. Eigen-Oscillation Spectrum and Active Dimension (O(N))}}\label{eigen-oscillation-spectrum-and-active-dimension-on}

The nonzero eigenvalues of the bipartite structure
\(W=\begin{pmatrix}0&B\\B^{\mathsf T}&0\end{pmatrix}\) are
\(\pm\sqrt{\operatorname{eig}(BB^{\mathsf T})}\). Indexing cross-parity
edges by an \(a\times b\) lattice of (even vertex)×(odd vertex),

\[BB^{\mathsf T}=(a-2)I_a\otimes J_b+(b-2)J_a\otimes I_b+2J_a\otimes J_b,\]

with eigenvalues \(4ab-2(a+b)\) (\(\times1\)), \(b(a-2)\)
(\(\times(a-1)\)), \(a(b-2)\) (\(\times(b-1)\)), and the rest \(0\).
Substituting \(a,b\) (even \(N\): \(a=b=N/2\); odd \(N\):
\(a=\tfrac{N+1}2,b=\tfrac{N-1}2\)):

\begin{itemize}
\tightlist
\item
  \textbf{Even \(N\ge6\)}: \(\lambda_{\max}^2=N(N-2)\) (multiplicity 1),
  \(\lambda_1^2=N(N-4)/4\) (multiplicity \(N-2\)).
\item
  \textbf{Odd \(N\ge7\)}: \(\lambda_{\max}^2=N^2-2N-1\) (multiplicity
  1), \(\lambda_1^2=(N-5)(N+1)/4\) (multiplicity \((N-3)/2\)),
  \(\lambda_2^2=(N-1)(N-3)/4\) (multiplicity \((N-1)/2\)).
\end{itemize}

Numerically matched for all \(N\) (even \(N\ge6\) / odd \(N\ge7\)).
\(N=3,4,5\) are degenerate and some eigenvalue families vanish (\(N=4\):
\(\lambda_1^2=0\); \(N=5\): \((N-5)(N+1)/4=0\)).

\textbf{Active dimension \(O(N)\).} The number of nonzero eigenvalues is
\(a+b-1=N-1\) (with \(\pm\) pairs, \(2(N-1)\)). Hence

\[\operatorname{rank}K=\operatorname{rank}W=2(N-1),\qquad \dim\ker=M-2(N-1)=\frac{(N-1)(N-4)}2\quad(\text{even }N\ge6,\ \text{odd }N\ge7).\]

Of the \(M=N(N-1)/2=O(N^2)\)-dimensional relational-wave space, only
\(O(N)\) dimensions are active under the 90-degree floor's generator;
the rest fall into \(\ker K\). Small \(N\) are exceptions due to
degeneracy: \(N=3\) (rank 2, kernel 1), \(N=4\) (rank 2, kernel 4),
\(N=5\) (rank 6, kernel 4).

A simple integer-harmonic system (integer multiples of a fundamental)
occurs only for \(N=3,4\). For \(N\ge5\) the ratios are in general
square-root ratios. Comparing the raw eigenvalues \(\lambda\) across
\(N\), there are 65 pairs of machine-precision integer ratios of
\(2\)--\(8\), and with the phase step included, \(q=\lambda/N\) detects
11 pairs (\texttt{固有振動数解析\_90度系列\_N3\_N40\_20260910/}). This
is consistent with the floor being a single eigenmode that co-winds a
single carrier (Paper 4: even after factorization, the frequency ratio
on the floor is 1:1).

\subsection{7. Relation to the Non-90-Degree Theoretical
Floor}\label{relation-to-the-non-90-degree-theoretical-floor}

{\def\LTcaptype{none} % do not increment counter
\begin{longtable}[]{@{}
  >{\raggedright\arraybackslash}p{(\linewidth - 8\tabcolsep) * \real{0.2000}}
  >{\raggedright\arraybackslash}p{(\linewidth - 8\tabcolsep) * \real{0.2000}}
  >{\raggedright\arraybackslash}p{(\linewidth - 8\tabcolsep) * \real{0.2000}}
  >{\raggedright\arraybackslash}p{(\linewidth - 8\tabcolsep) * \real{0.2000}}
  >{\raggedright\arraybackslash}p{(\linewidth - 8\tabcolsep) * \real{0.2000}}@{}}
\toprule\noalign{}
\begin{minipage}[b]{\linewidth}\raggedright
\end{minipage} & \begin{minipage}[b]{\linewidth}\raggedright
Phase
\end{minipage} & \begin{minipage}[b]{\linewidth}\raggedright
Generator
\end{minipage} & \begin{minipage}[b]{\linewidth}\raggedright
Solution method
\end{minipage} & \begin{minipage}[b]{\linewidth}\raggedright
Symmetry
\end{minipage} \\
\midrule\noalign{}
\endhead
\bottomrule\noalign{}
\endlastfoot
Cyclic-Fourier floor (Paper 2's (C)) & \(2\pi(i+j)/N\) (N-valued) &
continuous & distance-class-reduced eigenvector & cyclic symmetry
(amplitude constant within a distance class) \\
90-degree integer-K floor (this paper) & \(0/90/180/270\) (4-valued) &
\textbf{integer matrix} & \(\mu=-\sigma_{\max}\) eigenvector & breaks
distance-class symmetry (amplitude per parity class) \\
\end{longtable}
}

The 90-degree integer-K floor is not distance-class symmetric (amplitude
per parity class: 2-valued for even \(N\), 3-valued for odd \(N\)), so
it cannot be built by distance-class reduction; the correct approach is
to solve it as the Perron mode of the bipartite \(W\) (= the
\(\sigma_{\max}\) eigenmode of the integer \(K\)). The two floors are
distinct members, but both satisfy the ignition qualification of Paper 2
(an unstable relative equilibrium).

\subsection{8. Claim Classification · Open
Items}\label{claim-classification-open-items}

\begin{itemize}
\tightlist
\item
  Classification: analytic exact solution (parity bipartite structure +
  reduced Perron problem) + symbolic (arbitrary \(N\), sympy) /
  numerical (\(N=3\)--\(40\)) verification. Verdict: retained.
\item
  \textbf{Resolved (given closed form / general proof in this version)}:
  (1) The edge assignment of the amplitude classes is fixed to be
  exactly the parity classes themselves (cross-parity = \(i+j\) odd /
  even-side internal = both \(i,j\) even / odd-side internal = both
  \(i,j\) odd) (§5). (2) The residual centroid for odd \(N\) is also in
  closed form (\(D=N^2-2N-1\)):
\end{itemize}

\[|\textstyle\sum z|=\begin{cases}0&N\equiv0\ (\mathrm{mod}\,4)\\[1mm] \sqrt2/N&N\equiv2\ (\mathrm{mod}\,4)\\[1mm] \sqrt{(N-1)/(2(N+1)D)}&N\equiv1\ (\mathrm{mod}\,4)\\[1mm] \sqrt{(N+1)/(2(N-1)D)}&N\equiv3\ (\mathrm{mod}\,4)\end{cases}\]

Matches the measured values for all \(N=3\)--\(40\) (error
\(\le10^{-9}\); e.g.~\(N=5\): 0.154303, \(N=6\): 0.235702, \(N=7\):
0.140028). (3) Square-closure \(\sum z^2=0\) is a consequence of the
bipartite eigenmode, a general theorem holding not only for the Perron
mode but for every nonzero eigenmode (§2). (4) \textbf{The connectivity
of the bipartite \(W\) is proven in general for arbitrary \(N\ge3\)
(recovering the open item of Paper 1)}: even vertices \(U\)
(\(|U|=\lceil N/2\rceil\ge2\)), odd vertices \(V\). Cross-parity edges
\((u,v),(u',v')\) are connected by \((u,v)\)--\((u,u')\)--\((u',v')\)
(\((u,u')\) is even-side internal) if \(u\neq u'\), and by
\((u,v)\)--\((v,v')\)--\((u,v')\) if \(u=u'\) (\(v\neq v'\),
\(|V|\ge2\)). A same-parity edge \((u,u')\) connects to its component
via \((u,u')\)--\((u,v)\) (any \(v\in V\)). Hence for \(N\ge3\) it is
connected = irreducible, and Perron--Frobenius is applicable for
arbitrary \(N\ge3\) (numerically confirmed for \(N=3\)--\(40\) that the
number of connected components \(=1\)). (5) Enumeration of the
sign-gauge family: independently for each edge \(k_e\to k_e+2t_e\)
(\(S=\operatorname{diag}((-1)^{t_e})\), \(K\to SKS\) with \(W\)
invariant) gives \(2^M\) possibilities, and identifying the all-edge
simultaneous flip \(S=-I\) (global phase \(\pi\)) yields a relative sign
gauge of \(2^{M-1}\). - Open: within the scope of this paper no
structural open items remain. Classification of floors by equivalences
other than the sign gauge, such as vertex permutations (organizing the
geometric moduli), is outside the scope of this paper.

\subsection{Related Work (Structural Agreement · Not the Derivation
Source)}\label{related-work-structural-agreement-not-the-derivation-source}

\begin{itemize}
\tightlist
\item
  Spectra of line graphs (Cvetković et al., algebraic graph theory): the
  eigenvalues of \(A=L(K_N)\) and the integer antisymmetric operator on
  it.
\item
  Relative equilibria (Marsden \& Ratiu): the floor = a rigid-rotation
  relative equilibrium.
\end{itemize}

\subsection{References}\label{references}

\begin{enumerate}
\def\labelenumi{\arabic{enumi}.}
\tightlist
\item
  Noriaki Kihara, ``The Mechanism of Inflationary Rapid Expansion in
  Self-Consistent Relational-Wave Closed Systems,'' Concept DOI
  10.5281/zenodo.22112008.
\item
  D. Cvetković, P. Rowlinson, S. Simić, \emph{An Introduction to the
  Theory of Graph Spectra}, Cambridge (2010).
\end{enumerate}

\subsection{Reproduction}\label{reproduction}

The generator, verification, and data are in
\texttt{90度理論床\_整数K解析解\_全N\_20260909/}
(\texttt{make\_90deg\_floor\_analytic\_v1.py} = generator,
\texttt{closed\_form\_90deg\_floor\_verify\_v1.py} = closed-form sympy
exact verification,
\texttt{parents\_90deg\_floor\_analytic/parent\_90deg\_floor\_N\{00003..00040\}\_analytic.npz}
= 38 floors, \texttt{summary\_90deg\_floor\_closed\_form.csv},
\texttt{manifest\_90deg\_floor\_analytic.json},
\texttt{fig\_90deg\_floor\_complex\_planes\_analytic.png},
\texttt{分析\_90度理論床\_整数K解析解\_全N\_20260909.md},
\texttt{run\_all.sh}, \texttt{SHA256SUMS.txt}). The verification of the
arbitrary-\(N\) analytic reduction via the parity bipartite structure
(§2 · §3 · §5 · §6 · §8) is in
\texttt{同/parity二部構造\_解析的証明\_20260912/}
(\texttt{verify\_parity\_bipartite\_analytic\_20260912.py} = read-only
verification importing the canonical adjacency/one\_step with SHA check
+ sympy symbolic, \texttt{results/parity\_bipartite\_summary.csv},
\texttt{results/実行ログ\_20260912.log},
\texttt{分析\_parity二部構造\_解析的一本化\_20260912.md},
\texttt{README.md}, \texttt{run\_all.sh}, \texttt{SHA256SUMS.txt}). The
eigen-oscillation spectrum is in
\texttt{同/固有振動数解析\_90度系列\_N3\_N40\_20260910/}. The
equivalence of the complex self-consistent problem ≡ the real symmetric
\(W\) (\(D^{-1}KD=iW\)) and Perron--Frobenius are in Paper 1.
adjacency/one\_step are extracted from the canonical source
\texttt{run\_N3\_N40\_stage123\_v1.py} (SHA256
\texttt{1abf2353fee2e4f56f05e7a6f149fd086885136beb61ab571b48a56b09691567})
with no physics changed.

\end{document}
